% ===================================================================== MODEL
\section{System and threat model}
\label{sec:model}

\subsection{Network, committee, and adversary}
\sage{} targets monolithic permissioned chains with deterministic finality and a fixed migration
committee; probabilistic-finality sources, concurrent membership change, and sharded state are outside
the model. We assume partial synchrony~\cite{dls}. Before an unknown global stabilization time (GST), message
delays are arbitrary; after GST they are bounded by $\Delta$. Protocol messages travel over
authenticated channels. The analysis fixes one committee
$\VM=\mathcal V_1=\mathcal V_2$ of size $n$, with no membership change during $[\hd,\hr)$. A single
corruption set contains at most $f\le\min(f_1,f_2)$ validators, where $f_1$ and $f_2$ are the fault
thresholds of the source and target engines. The adversary schedules crashes and partitions, delays
messages before GST, and cannot forge signatures or find hash collisions. Distinct source and target
committees would require a joint-quorum rule and an explicit key mapping; our theorems do not cover
that setting.

\paragraph{Cross-domain Byzantine capability}
A Byzantine validator may act in \emph{both} protocol domains during partial activation. It may sign
source-engine and target-engine messages in the same execution, release a fence share that it never
durably installed, and authorize source finalization at or above the boundary while it also votes in
the active target. None of this requires equivocating between two legacy blocks, so a source-side
non-equivocation assumption does not exclude it. In the fence-only counterexample of
\cref{cor:cross-split}, this capability lets the residual source and target quorums overlap only in
the corrupted set. The admitted locked profile instead excludes that residual source quorum through
\cref{rule:lock,lem:fence-live}; the separate partial-delivery model (\ref{app:spec}) checks the
fence-only contract at bounded sizes without weakening its Byzantine capability.

\subsection{Replicated state and the engine interface}
Let the consensus-independent replicated state at height $h$ be
\begin{equation}
  X_h=(A_h,P_h),\qquad r_h=\mathrm{root}(X_h),\qquad X_h=\delta_{\mathrm{exec}}(X_{h-1},B_h),
  \label{eq:state}
\end{equation}
where $A_h$ is the execution-layer account and storage state, $P_h$ is platform state (including the
validator registry and the migration schedule), and $r_h$ is the post-execution root committed in the
header of $B_h$. Each engine $j$ also keeps local consensus metadata $C^{(j)}_h$ (views, locks,
certificates). This metadata is deliberately excluded from $r_h$, because heterogeneous engines encode
it differently. \Cref{tab:notation} lists the notation.

\begin{definition}[Consensus engine]
\label{def:engine}
A consensus engine is a tuple $\mathcal C=(\mathsf Q,\mathsf{Final},\mathsf{Root},\mathsf{Init})$.
Here $\mathsf Q$ is a finalization quorum rule, and $\mathsf{Final}$ is a single-finalizer predicate
that finalizes at most one block per height once a $\mathsf Q$-quorum endorses it. $\mathsf{Root}$ is
a deterministic state-root function, and $\mathsf{Init}(r,h)$ initializes $C^{(j)}_h$ from a root and
a height. The engine is \emph{intersection-safe} if its threshold $q$ satisfies $2q-n>f$. The minimum
such integer is $\lfloor (n+f)/2\rfloor+1$, which equals $2f+1$ only when $n=3f+1$.
\end{definition}

\Cref{def:engine} is an assume--guarantee interface. Real BFT safety also depends on durable locks,
view change, and crash recovery, which the tuple does not reconstruct. \sage{} establishes boundary
properties \emph{conditional} on each engine conforming to its interface. $\mathsf{Final}$ requires
\emph{deterministic} finality. This excludes longest-chain sources, whose ``finality'' is a
confirmation-depth confidence; migrating from such a source would first require a rule for selecting
and fencing a deterministic prefix. Applicability further requires deterministic state execution, a
shadow-validation adapter, externally bootstrappable target metadata, durable Rule~\ref{rule:lock}
and generation fencing, and a registered fixed committee. The evaluated PoA$\to$HotStuff
network path is partial; Raft$\to$HotStuff is exercised only as a component path
(\cref{tab:status}).

\Cref{tab:adapter} makes adapter obligations explicit. A new target must discharge native
finality, shadow, bootstrap, and liveness requirements; the source and controller must also meet
retirement, recovery, and terminal requirements. Naming an engine proves neither integration
nor compatibility. Its finality must remain safe across views and restart, not merely satisfy
a one-view threshold.

\begin{table*}[t]
\centering
\caption{Assume--guarantee obligations for a concrete engine adapter. Listed evidence is required,
not a claim of a full-protocol run or verified Mysticeti, Bullshark, or HotStuff-1 adapters.}
\label{tab:adapter}
\small
\renewcommand{\arraystretch}{1.12}
\begin{tabularx}{\textwidth}{@{}L{0.18\textwidth}YL{0.28\textwidth}@{}}
\toprule
\textbf{Obligation} & \textbf{Required guarantee} & \textbf{Evidence required of an instantiation} \\
\midrule
Finality and fault profile & Deterministic finality, authenticated committee identities, quorum policy and admitted corruption set & Engine safety argument including cross-view locking; fault budget and voting-set checks \\
State and shadow & Same deterministic transition and root on the same ordered inputs; no shadow finalization or external side effects & Root-equivalence and rejection tests using real executed state \\
Bootstrap & Legacy anchor distinguished from a native target QC; first native certificate obtained from a target quorum & First target proposal/vote/commit trace and rejection of mismatched anchors \\
Durable source retirement & Persist attestation identity and prohibition atomically before release; serialize every source authorization path; retain later fences & Crash-before/after-write, buffered-send, alternate-handler and post-restart source-action traces \\
Recovery and activation & Reload Rule~\ref{rule:lock} before source use; reload manifest, generation and certificate records before target authority; fail closed on ambiguity & Reopen/corruption tests and first post-restart signed actions \\
Conditional liveness & Bounded certificate delivery and durable work; target-specific first-finalization bound after initialization & Engine-specific pacemaker argument and wall-clock execution, not a generic $2\Delta$ assertion \\
Terminal/client semantics & Abort and seal are exclusive; recovery never revives $g_1$; effective finality retains the original signed receipt & Certified abort/seal race and recovery traces; contextual replay and client proof checks \\
\bottomrule
\end{tabularx}
\end{table*}

\begin{definition}[Target bootstrap]
\label{def:bootstrap}
At the certified transition height $\hs=\hb+1\ge\hc$, the target initializes
$C^{(2)}_{\hs}\gets\Ctgt.\mathsf{Init}(r_{\hs-1},\hs)$. It adopts the legacy-finalized root as its
genesis anchor and treats the manifest's bootstrap certificate as an authority-transfer anchor, not as
an ordinary target quorum certificate. In HotStuff, the highest quorum certificate (QC) and lock start empty, and the first
target-native QC must be formed by a real target quorum. $A$ and $P$ carry over unchanged.
\end{definition}

A same-engine reconfiguration may also require state transfer and initialization. A heterogeneous
handoff additionally must translate the certified boundary into the target engine's native metadata;
the target has no native consensus history on the legacy chain.

\begin{table}[t]
\centering
\caption{Notation.}
\label{tab:notation}
\small
\renewcommand{\arraystretch}{1.12}
\rowcolors{2}{rowtint}{white}
\begin{tabular}{@{}>{$}l<{$}L{0.66\columnwidth}@{}}
\toprule
\textbf{Symbol} & \textbf{Meaning} \\
\midrule
\VM,\ n,\ f & fixed migration committee, its size, Byzantine bound ($n\ge3f+1$) \\
\Csrc,\ \Ctgt & source (legacy) and target consensus engines \\
g_1,\ g_1',\ g_2 & source, fresh recovery-source, and target generations \\
q,\ q_1,\ q_2,\ q_F & certificate, source, target, and fence quorum sizes \\
X_h,\ r_h,\ C^{(j)}_h & replicated state, its root, engine-local metadata \\
\hd,\ \hc,\ \hs,\ \hr & dual-run start, earliest cutover, certified cutover, rollback deadline \\
\sigma_h,\ \tilde r_h,\ \kappa & shadow verdict, shadow root, stability window \\
\beta,\ \mathcal M_{\mathrm B} & boundary tuple and complete manifest body \\
\pi_{\mathrm B},\ \pi_{\mathrm M} & \CutCert{} and \ManCert{} \\
\Auth(h),\ \tau & authoritative engine at $h$; target progress timeout \\
\bottomrule
\end{tabular}
\end{table}

\subsection{Fault thresholds of the two engines}
\label{subsec:thresholds}
Under its non-equivocation assumption, the PoA source finalizes with an honest-majority quorum
$\mathsf Q_1=\lfloor n/2\rfloor+1$. The target uses $\mathsf Q_2=q_2$ with $2q_2-n>f_2$ and
$n\ge 3f_2+1$. We instantiate $q_2=n-f_2$, since $2(n-f_2)-n=n-2f_2>f_2$ exactly when $n>3f_2$. The
smaller admissible $\lfloor(n+f_2)/2\rfloor+1$ is used only in a labelled research control.

A round-robin PoA source tolerates $f_1<n/2$ crash faults, but an honest-majority rule does not stop
a Byzantine authority from signing two blocks at one height. We therefore admit two source models.
Under the \emph{crash-only} model, authority keys are trusted not to equivocate. Under the
\emph{Byzantine} model, the source must enforce non-equivocation per key (for example through attested
append-only key custody) or use a quorum with $2q_1-n>f_1$. The evaluated PoA source uses the
crash-only interpretation. SeeMoRe~\cite{seemore} mixes crash-only and Byzantine behaviour inside a
single protocol. \sage{} faces the mixture across an engine boundary.

During dual-run, safety is bounded by the weaker source engine. This bound predates the migration: a
chain running $\Csrc$ already tolerates at most $f_1$ faults. \sage{} adds no finalizing authority
during $[\hd,\hs)$. Its shadow validation adds a detection surface for invalid or semantically
incompatible execution, but it cannot detect a well-formed malicious block that both engines execute
identically.

\subsection{Threats and gating mechanisms}
\label{sec:threat}
\Cref{fig:threat} places five attack surfaces on the migration timeline and pairs each with the
mechanism that gates it; \cref{tab:crypto} lists the assumptions each mechanism relies on.
(T1)~\emph{Key custody}: every certificate relies on signatures that are existentially unforgeable
under chosen-message attack (EUF-CMA) and on at most $f$ corrupted
keys. (T2)~\emph{False readiness}: an $n-f$ \CutCert{} contains at least $f+1$ correct signers, and a
live node checks the certified boundary against its own committed tip. (T3)~\emph{Partition at
cutover}: a side below $n-f$ cannot form a new \CutCert. Under durable Rule~\ref{rule:lock} and
$q_1>2f$, a formed \CutCert{} already precludes a source quorum even at non-recipients. Without
that lock, the separate fence-only contract must exclude overlap confined to Byzantine validators
(\cref{cor:cross-split}); that corollary's rejected source policy is not an admitted locked
execution. (T4)~\emph{Manifest supply chain}: the \ManCert{}
authenticates the complete manifest, including tooling and replay-context hashes. Quorum agreement on
a tool hash does not make the tool safe, so deployment still needs an allowlist or reproducible
builds. (T5)~\emph{Progress griefing}: a local timeout restores the certified anchor without changing
committee authority. Only an $n-f$ \AbortCert{} returns authority, and only to a fresh generation
$\Csrc^{g_1'}$ with $g_1'>g_1$.

\begin{figure*}[t]
\centering
\includegraphics[width=\textwidth]{figures/fig_threat.pdf}
\caption{Migration attack surfaces (T1--T5) and their protocol gates (M1--M6). The target suffix
$[\hs,\hr)$ remains provisional until a \SealCert{} promotes it to absolute finality.}
\label{fig:threat}
\end{figure*}

\begin{table}[t]
\centering
\caption{Cryptographic and combinatorial assumptions, and the result each supports.}
\label{tab:crypto}
\small
\renewcommand{\arraystretch}{1.14}
\rowcolors{2}{rowtint}{white}
\begin{tabular}{@{}L{0.42\columnwidth}L{0.52\columnwidth}@{}}
\toprule
\textbf{Assumption} & \textbf{Supports} \\
\midrule
Collision-resistant $\mathrm{root}(\cdot)$, $\Hash$ & boundary and manifest binding (\cref{th:unforge}) \\
EUF-CMA signatures & authenticity of every certificate share \\
Chain-id and epoch binding & replay resistance (\cref{prop:replay}) \\
$n\ge3f+1$, one committee & certificate intersection (\cref{th:window,th:unique}) \\
Durable Rule~\ref{rule:lock}, $q=n-f$, $q_1>2f$ & locked-profile source exclusion (\cref{lem:fence-live}) \\
$q_F+q_1>n+f$, install-before-share & fence-only exclusion (\cref{th:authority-exclusion}); not independently needed under Rule~\ref{rule:lock} \\
\bottomrule
\end{tabular}
\end{table}
